\documentclass[reprint,english,superscriptaddress,preprintnumbers,amsmath,amssymb,aps,prd]{revtex4-1}

%%%%%%%%%%%%% Packages %%%%%%%%%%%%%

% general
\usepackage[utf8]{inputenc}


\usepackage{diagbox}

% math
\usepackage{mathtools}
\usepackage{amsfonts}
\usepackage{mathrsfs}
\usepackage{bbm}
\usepackage{slashed}
\usepackage{braket}
\usepackage{titlesec}

% graphics and colors
\usepackage{graphicx}
\usepackage{color}
\usepackage{array}
% \usepackage{ulem}

% floats
\usepackage{blindtext}
\usepackage{placeins}
\usepackage{booktabs}
\usepackage{makecell}
\usepackage{epstopdf}
%\usepackage{subcaption}
%\usepackage{floatrow}
\usepackage[caption=false]{subfig}
%\usepackage{dblfloatfix} % fix for bottom-placement of figure

% units and refs
\usepackage{xspace}
\usepackage{siunitx}
\usepackage{hyperref}
\usepackage[nameinlink]{cleveref}
\usepackage{appendix}


% other
\usepackage{xifthen}
\usepackage{xcolor}
\hypersetup{
	colorlinks,
	linkcolor={blue!75!black},%{red!75!black},
	citecolor={blue!75!black},
	urlcolor={blue!75!black}
}

%%%%%%%%%%%%% Options %%%%%%%%%%%%%

% set default value of \includegraphics
\setkeys{Gin}{width=0.48\textwidth}

% justification of figure captions
\captionsetup{justification=centerlast}

\graphicspath{{figures/}}


%%%%%%%%%%%%% Commands %%%%%%%%%%%%%

% typesetting
\newcommand{\eqnewline}{\nonumber\\[2ex]}
\newcommand{\myhfill}{\hfill\textcolor{white}{.}}
\newcommand*{\eg}{e.g.\@\xspace}
\newcommand*{\ie}{i.e.\@\xspace}
\newcommand*{\etc}{etc.\@\xspace}
\newcommand*{\cf}{cf.\@\xspace}
\newcommand{\overbar}[1]{\mkern 2mu\overline{\mkern-3mu#1\mkern-1.7mu}\mkern 1.7mu}



%%%%%%%%%%%%% Referencing %%%%%%%%%%%%%

\def\eq#1{\eqref{#1}}
%\def\eqs#1{(\ref{#1})}
\def\Eq#1{Eq.~\labelcref{#1}}
\def\Eqs#1{Eqs.~\labelcref{#1}}

\def\fig#1{Fig.~\labelcref{#1}}
\def\Fig#1{Fig.~\labelcref{#1}}
\def\figs#1#2{Figs.~\labelcref{#1} and \labelcref{#2}}
\def\Figs#1#2{Figs.~\labelcref{#1} and \labelcref{#2}}
\def\Tab#1{Tab.~\labelcref{#1}}
\def\tab#1{Tab.?~\labelcref{#1}}
\def\Sec#1{Sec.~\labelcref{#1}}
%\def\Secs#1#2{Sections~\ref{#1} and \ref{#2}}
\def\App#1{App.~\labelcref{#1}}


\newcommand{\tr}{{\text{tr}}}

%%%%%%%%%%%%% Math commands %%%%%%%%%%%%%
% symbols
\newcommand{\sumint}{\int\hspace{-4.8mm}\sum}
\newcommand{\imag}{\text{i}}
\newcommand{\tinytext}[1]{\text{\tiny{#1}}}
\renewcommand{\Re}{\operatorname{Re}}
\renewcommand{\Im}{\operatorname{Im}}
\newcommand{\sgn}[1]{\mathrm{sgn}\left(#1\right)}
\usepackage{multirow}
\usepackage{colortbl}
\definecolor{kugray5}{RGB}{224,224,224}
\newcommand{\PreserveBackslash}[1]{\let\temp=\\#1\let\\=\temp}
\newcolumntype{C}[1]{>{\PreserveBackslash\centering}p{#1}}
\newcolumntype{R}[1]{>{\PreserveBackslash\raggedleft}p{#1}}
\newcolumntype{L}[1]{>{\PreserveBackslash\raggedright}p{#1}}


%%%%%%%%%%%%% Comments %%%%%%%%%%%%%
\newcommand{\colwj}[1]{\textcolor{red}{#1}}
\newcommand{\colch}[1]{\textcolor{cyan}{#1}}
\newcommand{\collc}[1]{\textcolor{magenta}{#1}}
\newcommand{\colyy}[1]{\textcolor{green}{#1}}
\newcommand{\cJMP}[1]{\textcolor{blue}{#1}}




%%%%%%%%%%%%% Title and hypersetup %%%%%%%%%%%%%


\newcommand{\gettitle}{Non-thermal fixed point}

\newcommand{\getHeidelbergAffiliation}{\affiliation{Institut f{\"u}r Theoretische Physik, Universit{\"a}t Heidelberg, Philosophenweg 16, 69120 Heidelberg, Germany}}

\newcommand{\getEMMIAffiliation}{\affiliation{ExtreMe Matter Institute EMMI, GSI, Planckstr. 1, D-64291 Darmstadt, Germany}}

\hypersetup{
	pdftitle={\gettitle},
	pdfauthor={Chuang, Pawlowski},
	pdfkeywords={functional methods}
		{Distribution Amplitudes}
		{functional renormalisation group},
	bookmarksopen=true,
	bookmarksopenlevel=2,
	bookmarksnumbered=true
}


\begin{document}


\title{\gettitle}

\author{Chuang Huang}
\email{huang@thphys.uni-heidelberg.de}
\getHeidelbergAffiliation

\author{Jan M. Pawlowski}
%\email{j.pawlowski@thphys.uni-heidelberg.de}
\getHeidelbergAffiliation\getEMMIAffiliation


%\pacs{Valid PACS appear here}% PACS, the Physics and Astronomy
\pacs{11.30.Rd, %Chiral symmetries
	%11.10.Wx, %Finite-temperature field theory
	12.38.Aw, %General properties of QCD (dynamics, confinement, etc.)
	05.10.Cc, %renormalisation group methods
	12.38.Mh,  %Quark-gluon plasma
	12.38.Gc %Lattice QCD calculations
}                             % Classification Scheme.
%\keywords{Suggested keywords}%Use showkeys class option if keyword
%display desired


%%%%%%%%%%%%%%%%%%%%%%%%%%%%%%%%%%%%%%%%%%%%%
\newpage


%\appendix
\renewcommand{\thesubsection}{{S.\arabic{subsection}}}
\setcounter{section}{0}
\titleformat*{\section}{\centering \Large \bfseries}

\onecolumngrid

	

%%%%%%%%%%%%%%%%%%%%%%%%%%%%%%%
\section*{Non-equilibrium real-time O(N) theory}



%%%%%%%%%%%%%%%%%%%%%%%%%%%%%%%%%%%%%%%%%%%%%%%%%%%%%%%%%%

\subsection{Effective action and correlation functions}
\label{sec:action}

With the definition of the fields,
\begin{align}
\phi_{\pm}=\frac{1}{\sqrt{2}}(\phi_c\pm\phi_q),\quad 
\phi_{c}=(\phi_{+}+\phi_{-})/\sqrt{2},\quad \phi_{q}=(\phi_{+}-\phi_{-})/\sqrt{2},
\end{align}
the effective action is given by
\begin{align}
\Gamma[\phi_c,\phi_q]=&\int_{p}\,\left\{\phi_{q}(-p)P^{R}(p)\phi_{c}(p)+\phi_{c}(-p)P^{A}(p)\phi_{q}(p)+\frac{1}{2}\phi_{q}(-p)\left
(\mathrm{i}D^{K}(p)\right)\phi_{q}(p)\right\}\nonumber
\\[2ex]
&-\int_{p_1,p_2,p_3,p_4}\left\{\frac{1}{6}\left(\phi_{q}(p_1) \phi_{c}(p_2)\right)\left(\lambda_{cl}\cdot \phi_{c}(p_3) \phi_{c}(p_4)+\lambda_{qu}\cdot \phi_{q}(p_3)\phi_{q}(p_4)\right)\right\}\delta(p_1+p_2+p_3+p_4)\nonumber\\[2ex]
&-\int_{p_1,p_2,p_3,p_4}\left\{\frac{\mathrm{i}}{2}\lambda_{an}\cdot \phi_{q}(p_1) \phi_{c}(p_2)\phi_{q}(p_3)\phi_{c}(p_4)\right\}\delta(p_1+p_2+p_3+p_4)\,.
\end{align}
Based on the effective action, we derive the propagators and the vertices. Firstly, the two-point correlation function matrix is given by
\begin{align}
    \Gamma^{(2)}=
\begin{pmatrix}
  0 & \Gamma^{cq} \\
  \Gamma^{qc} & \Gamma^{qq}
\end{pmatrix}=
\begin{pmatrix}
  0 & P^{A} \\
  P^{R} & \mathrm{i}D^{K}
\end{pmatrix}\,,
\end{align}
where $P^{R}=\left(P^{A}\right)^{*}$. Therefore, the propagator reads,
\begin{align}
    G=\left(\Gamma^{(2)}\right)^{-1}=\begin{pmatrix}
  G^{K} & G^{R} \\
  G^{A} & 0
\end{pmatrix}\,,
\end{align}
with
\begin{align}
    G^{R}=\left(P^{R}\right)^{-1},\quad G^{A}=\left(P^{A}\right)^{-1},\quad \text{and}\quad
    G^{K}=-\mathrm{i}\,G^{R}D^{K}G^{A}\,.
\end{align}
In the equilibrium, ultraviolet, and vacuum limit, the correlation functions are given by
\begin{align}
    P^{R}_{ij}(p)&=\delta_{ij}\left\{p_{0}^{2}-\boldsymbol{p}^{2}-m^{2}+\mathrm{i}\epsilon\,\mathrm{sgn}(p_{0})\right\}\,,\nonumber\\[2ex] 
    P^{A}_{ij}(p)&=\delta_{ij}\left\{p_{0}^{2}-\boldsymbol{p}^{2}-m^{2}-\mathrm{i}\epsilon\,\mathrm{sgn}(p_{0})\right\}\,
\end{align}
and 
\begin{align}
    \mathrm{i}D^{K}_{ij}=\delta_{ij}\left\{2\mathrm{i}\epsilon\,\mathrm{sgn}(p_{0})\right\}\,.
\end{align}

More importantly, for the vertex, we employ the $s,t,u$-channel approximation, where the details are in \Cref{tab:stu}. Each channel dressing depends on the two endpoint times of that channel and on a single channel momentum. The two-point flows derived below never call a vertex that depends on two independent channel momenta simultaneously. Accordingly, we close the vertex sector on four independent dressings, listed in \Cref{tab:dressings}: a single classical dressing $\lambda_{cl}$, a single quantum dressing $\lambda_{qu}$, and two anomaly dressings $\lambda_{an}^{s}$ and $\lambda_{an}^{t}$.

The identification $\lambda_{cl}^{s}=\lambda_{cl}^{t}=\lambda_{cl}^{u}\equiv\lambda_{cl}$ (and likewise for $\lambda_{qu}$) follows from the $1+3$ leg structure of the $qccc$ and $qqqc$ vertices: the three index pairings are equivalent under permutations of the three identical legs, so the three channels obey the same flow with the same initial condition. By contrast, the $qqcc$ vertex has a $2+2$ leg structure. Its $s$-channel tensor $\delta_{ij}\delta_{kl}$ (amplitude-dependent noise) and $t/u$-channel tensors $\delta_{ik}\delta_{jl}+\delta_{il}\delta_{jk}$ (orientation-dependent noise) are independent $\mathrm{O}(N)$ invariants, so we keep
\begin{align}
\lambda_{an}^{s}\neq\lambda_{an}^{t}=\lambda_{an}^{u}\,.
\label{eq:an-stu}
\end{align}
At tree level $\lambda_{cl}=\lambda_{qu}$ and $\lambda_{an}=0$; the latter remains zero in equilibrium. One may optionally identify $\lambda_{cl}=\lambda_{qu}$ along the flow.

According to the definition, specifically, the classical vertex is given by 
\begin{align}
    \Gamma_{q_{i}c_{j}c_{k}c_{l}}^{(4)}=-\frac{1}{3}&\Big\{\lambda_{cl}\left(t_i,t_k,\left(\boldsymbol{p}_{1}+\boldsymbol{p}_{2}\right)^{2}\right)\delta_{ij}\delta_{kl}\delta(t_i-t_j)\delta(t_k-t_l)\nonumber\\[2ex]
    &+\lambda_{cl}\left(t_i,t_j,\left(\boldsymbol{p}_{1}+\boldsymbol{p}_{3}\right)^{2}\right)\delta_{ik}\delta_{jl}\delta(t_i-t_k)\delta(t_j-t_l)\nonumber\\[2ex]
    &+\lambda_{cl}\left(t_i,t_j,\left(\boldsymbol{p}_{1}+\boldsymbol{p}_{4}\right)^{2}\right)\delta_{il}\delta_{jk}\delta(t_i-t_l)\delta(t_j-t_k)\Big\}\,.\label{eq:classical-vertex}
\end{align}
Similarly, the quantum vertex is
\begin{align}
    \Gamma_{q_{i}q_{j}q_{k}c_{l}}^{(4)}=-\frac{1}{3}&\Big\{\lambda_{qu}\left(t_i,t_k,\left(\boldsymbol{p}_{1}+\boldsymbol{p}_{2}\right)^{2}\right)\delta_{ij}\delta_{kl}\delta(t_i-t_j)\delta(t_k-t_l)\nonumber\\[2ex]
    &+\lambda_{qu}\left(t_i,t_j,\left(\boldsymbol{p}_{1}+\boldsymbol{p}_{3}\right)^{2}\right)\delta_{ik}\delta_{jl}\delta(t_i-t_k)\delta(t_j-t_l)\nonumber\\[2ex]
    &+\lambda_{qu}\left(t_i,t_j,\left(\boldsymbol{p}_{1}+\boldsymbol{p}_{4}\right)^{2}\right)\delta_{il}\delta_{jk}\delta(t_i-t_l)\delta(t_j-t_k)\Big\}\,.\label{eq:quantum-vertex}
\end{align}
and the anomaly vertex is
\begin{align}
    \Gamma_{q_{i}q_{j}c_{k}c_{l}}^{(4)}=-\mathrm{i}&\Big\{\lambda_{an}^{s}\left(t_i,t_k,\left(\boldsymbol{p}_{1}+\boldsymbol{p}_{2}\right)^{2}\right)\delta_{ij}\delta_{kl}\delta(t_i-t_j)\delta(t_k-t_l)\nonumber\\[2ex]
    &+\lambda_{an}^{t}\left(t_i,t_j,\left(\boldsymbol{p}_{1}+\boldsymbol{p}_{3}\right)^{2}\right)\delta_{ik}\delta_{jl}\delta(t_i-t_k)\delta(t_j-t_l)\nonumber\\[2ex]
    &+\lambda_{an}^{t}\left(t_i,t_j,\left(\boldsymbol{p}_{1}+\boldsymbol{p}_{4}\right)^{2}\right)\delta_{il}\delta_{jk}\delta(t_i-t_l)\delta(t_j-t_k)\Big\}\,.\label{eq:anomaly-vertex}
\end{align}
where the anomaly vertex vanishes in the equilibrium case.

%
%%%%%%%%%%%%%%%%%%%%%%%%%%%%%
\begin{table*}[t]
    \begin{center}
    \centering
    \begin{tabular}{c|c|c|c}
      \hline\hline & & &   \\[-2ex]   
      Channel & Index pair & Time & Moment   \\[1ex]
      \hline & & &   \\[-2ex]
      \quad $s$-channel \quad\quad  & \quad $(i,j)|(k,l)$ \quad\quad  & \quad $\delta(t_i-t_j)\delta(t_k-t_l)$ \quad\quad  & \quad $s=\left(\boldsymbol{p}_{1}+\boldsymbol{p}_{2}\right)^{2}$ \quad\quad \\[1ex]
      \quad $t$-channel \quad\quad  & \quad $(i,k)|(j,l)$ \quad\quad  & \quad $\delta(t_i-t_k)\delta(t_j-t_l)$ \quad\quad  & \quad $t=\left(\boldsymbol{p}_{1}+\boldsymbol{p}_{3}\right)^{2}$ \quad\quad \\[1ex]
      \quad $u$-channel \quad\quad  & \quad $(i,l)|(j,k)$ \quad\quad  & \quad $\delta(t_i-t_l)\delta(t_j-t_k)$ \quad\quad  & \quad $u=\left(\boldsymbol{p}_{1}+\boldsymbol{p}_{4}\right)^{2}$ \quad\quad \\[1ex]
      \hline\hline
    \end{tabular}
    \caption{Definition of the $s,t,u$ channels for four-point functions. Each channel dressing depends on the two endpoint times of that channel and on the corresponding single channel momentum.}
    \label{tab:stu}
    \end{center}\vspace{-0.5cm}
\end{table*}
%%%%%%%%%%%%%%%%%%%%%%%%%%%%%
%
%%%%%%%%%%%%%%%%%%%%%%%%%%%%%
\begin{table*}[t]
    \begin{center}
    \centering
    \begin{tabular}{c|c|c|c}
      \hline\hline & & &   \\[-2ex]
      Dressing & Vertex & Channel content & Arguments \\[1ex]
      \hline & & &   \\[-2ex]
      $\lambda_{cl}(t_a,t_b;\boldsymbol{P}^{2})$ & $qccc$ & $\lambda_{cl}^{s}=\lambda_{cl}^{t}=\lambda_{cl}^{u}$ & two times, one momentum \\[1ex]
      $\lambda_{qu}(t_a,t_b;\boldsymbol{P}^{2})$ & $qqqc$ & $\lambda_{qu}^{s}=\lambda_{qu}^{t}=\lambda_{qu}^{u}$ & two times, one momentum \\[1ex]
      $\lambda_{an}^{s}(t_a,t_b;\boldsymbol{P}^{2})$ & $qqcc$ & $s$-channel ($\kappa_{1}$) & two times, one momentum \\[1ex]
      $\lambda_{an}^{t}(t_a,t_b;\boldsymbol{P}^{2})$ & $qqcc$ & $t=u$-channel ($\kappa_{2}$) & two times, one momentum \\[1ex]
      \hline\hline
    \end{tabular}
    \caption{Independent vertex dressings after the channel truncation. The two-point flows only evaluate these dressings at vanishing channel momentum (local $s$-type insertions) or at a single off-shell channel momentum $(\boldsymbol{p}\mp\boldsymbol{q})^{2}$ (memory insertions).}
    \label{tab:dressings}
    \end{center}\vspace{-0.5cm}
\end{table*}
%%%%%%%%%%%%%%%%%%%%%%%%%%%%%
%%%%%%%%%%%%%%%%%%%%%%%%%%%%%%%%%%%%%%%%%%%%%%%%%%%%%%%

\subsection{The flow equations for the two-point functions}
\label{sec:2pointflow}

For the quantum-classical correlation function, the flow equation is given by
\begin{align}
    \partial_{\tau}\Gamma_{q_i,c_j}^{(2)}(t_i,t_j;\boldsymbol{p})=&\frac{\mathrm{i}}{2}\tilde{\partial}_{\tau}\int_{\boldsymbol{q}}\int_{t_k}\int_{t_l}\mathrm{tr}\frac{1}{N}\delta_{ij}\nonumber\\[2ex]&\hspace{-2cm}\times\left\{\Gamma_{q_{i}c_{j}c_{k}c_{l}}^{(4)}(t_i,t_j,t_k,t_l;\boldsymbol{p},-\boldsymbol{p},\boldsymbol{-q},\boldsymbol{q})G^{K}_{kl}(t_k,t_l;\boldsymbol{q})\right.\label{eq:qc-1}\\[2ex]
    &\hspace{-1.5cm}+\Gamma_{q_{i}c_{j}q_{k}c_{l}}^{(4)}(t_i,t_j,t_k,t_l;\boldsymbol{p},-\boldsymbol{p},\boldsymbol{-q},\boldsymbol{q})G^{A}_{kl}(t_k,t_l;\boldsymbol{q})\label{eq:qc-2}\\[2ex]
    &\left.\hspace{-1.5cm}+\Gamma_{q_{i}c_{j}c_{k}q_{l}}^{(4)}(t_i,t_j,t_k,t_l;\boldsymbol{p},-\boldsymbol{p},\boldsymbol{-q},\boldsymbol{q})G^{R}_{kl}(t_k,t_l;\boldsymbol{q})\right\}\,,\label{eq:qc-3}
\end{align}
which includes the contribution from three diagrams. Here $\tau$ is the RG-time. Using \Cref{eq:classical-vertex} together with $G^{K}(t_j,t_i;-\boldsymbol{q})=G^{K}(t_i,t_j;\boldsymbol{q})$, \Cref{eq:qc-1} reduces to
\begin{align}
    \partial_{\tau}\Gamma_{q_i,c_j}^{(2),\mathrm{I}}(t_i,t_j;\boldsymbol{p})=&-\frac{\mathrm{i}}{6}\Bigg(N\delta(t_i-t_j)\int_{\boldsymbol{q}}\int_{t_k}\lambda_{cl}(t_i,t_k;0)\partial_{\tau}G^{K}(t_k,t_k;\boldsymbol{q})\nonumber\\[2ex]
    &\hspace{0.8cm}+2\int_{\boldsymbol{q}}\lambda_{cl}(t_i,t_j;(\boldsymbol{p}-\boldsymbol{q})^2)\partial_{\tau}G^{K}(t_i,t_j;\boldsymbol{q})\Bigg)\,.\label{eq:qc-I}
\end{align}
The first line is the local $s$-channel insertion, which renormalizes the mass. The second line is the $t+u$ memory insertion and retains $t_i-t_j\neq 0$. To match the definition \Cref{eq:anomaly-vertex}, \Cref{eq:qc-2} can be rewritten as
\begin{align}
    \partial_{\tau}\Gamma_{q_i,c_j}^{(2),\mathrm{II}}(t_i,t_j;\boldsymbol{p})&=\frac{\mathrm{i}}{2}\tilde{\partial}_{\tau}\int_{\boldsymbol{q}}\int_{t_k}\int_{t_l}\mathrm{tr}\frac{1}{N}\delta_{ij}\Gamma_{q_{i}c_{j}q_{k}c_{l}}^{(4)}(t_i,t_j,t_k,t_l;\boldsymbol{p},-\boldsymbol{p},\boldsymbol{-q},\boldsymbol{q})G^{A}_{kl}(t_k,t_l;\boldsymbol{q})\nonumber\\[2ex]
    &=\frac{\mathrm{i}}{2}\tilde{\partial}_{\tau}\int_{\boldsymbol{q}}\int_{t_k}\int_{t_l}\mathrm{tr}\frac{1}{N}\delta_{ij}\Gamma_{q_{i}q_{k}c_{j}c_{l}}^{(4)}(t_i,t_k,t_j,t_l;\boldsymbol{p},-\boldsymbol{q},\boldsymbol{-p},\boldsymbol{q})G^{A}_{kl}(t_k,t_l;\boldsymbol{q})\nonumber\\[2ex]
    &=\frac{\mathrm{1}}{2}\Bigg(N\delta(t_i-t_j)\int_{\boldsymbol{q}}\int_{t_k}\lambda_{an}^{t}(t_i,t_k;0)\partial_{\tau}G^{A}(t_k,t_k;\boldsymbol{q})\nonumber\\[2ex]
    &\hspace{0.5cm}+\int_{\boldsymbol{q}}\lambda_{an}^{s}(t_i,t_j;(\boldsymbol{p}-\boldsymbol{q})^2)\partial_{\tau}G^{A}(t_i,t_j;\boldsymbol{q})+\int_{\boldsymbol{q}}\lambda_{an}^{t}(t_i,t_j;(\boldsymbol{p}+\boldsymbol{q})^2)\partial_{\tau}G^{A}(t_j,t_i;\boldsymbol{q})\Bigg)\,.
\end{align}
Similarly, \Cref{eq:qc-3} can also be simplified. Using $G^{R}(-\boldsymbol{q})=G^{A}(\boldsymbol{q})$ and $\lambda_{an}^{t}=\lambda_{an}^{u}$, the retarded and advanced anomaly insertions combine. Finally, the flow equation of the quantum-classical two-point function reads
\begin{align}
    \partial_{\tau}\Gamma_{q_i,c_j}^{(2)}(t_i,t_j;\boldsymbol{p})=&-\frac{\mathrm{i}}{6}\Bigg(N\delta(t_i-t_j)\int_{\boldsymbol{q}}\int_{t_k}\lambda_{cl}(t_i,t_k;0)\partial_{\tau}G^{K}(t_k,t_k;\boldsymbol{q})\nonumber\\[2ex]
    &\hspace{0.5cm}+2\int_{\boldsymbol{q}}\lambda_{cl}(t_i,t_j;(\boldsymbol{p}-\boldsymbol{q})^2)\partial_{\tau}G^{K}(t_i,t_j;\boldsymbol{q})\Bigg)\nonumber\\[2ex]
    &+\frac{1}{2}N\delta(t_i-t_j)\int_{\boldsymbol{q}}\int_{t_k}\lambda_{an}^{t}(t_i,t_k;0)\partial_{\tau}\big(G^{R}+G^{A}\big)(t_k,t_k;\boldsymbol{q})\nonumber\\[2ex]
    &+\int_{\boldsymbol{q}}\Big[\lambda_{an}^{s}(t_i,t_j;(\boldsymbol{p}-\boldsymbol{q})^2)\partial_{\tau}G^{A}(t_i,t_j;\boldsymbol{q})+\lambda_{an}^{t}(t_i,t_j;(\boldsymbol{p}-\boldsymbol{q})^2)\partial_{\tau}G^{R}(t_i,t_j;\boldsymbol{q})\Big]\,.\label{eq:qc-final}
\end{align}
The two memory kernels $\lambda_{an}^{s}$ and $\lambda_{an}^{t}$ remain distinct: they multiply $G^{A}$ and $G^{R}$, respectively, and cannot be merged into a single $(G^{R}+G^{A})$ structure.

For the quantum-quantum correlation function, the same channel identifications yield
\begin{align}
    \partial_{\tau}\Gamma_{q_i,q_j}^{(2)}(t_i,t_j;\boldsymbol{p})=&\frac{1}{2}N\delta(t_i-t_j)\int_{\boldsymbol{q}}\int_{t_k}\lambda_{an}^{s}(t_i,t_k;0)\partial_{\tau}G^{K}(t_k,t_k;\boldsymbol{q})\nonumber\\[2ex]
    &+\int_{\boldsymbol{q}}\lambda_{an}^{t}(t_i,t_j;(\boldsymbol{p}-\boldsymbol{q})^2)\partial_{\tau}G^{K}(t_i,t_j;\boldsymbol{q})\nonumber\\[2ex]
    &-\frac{\mathrm{i}}{6}\Bigg(N\delta(t_i-t_j)\int_{\boldsymbol{q}}\int_{t_k}\lambda_{qu}(t_i,t_k;0)\partial_{\tau}\big(G^{R}+G^{A}\big)(t_k,t_k;\boldsymbol{q})\nonumber\\[2ex]
    &\hspace{0.5cm}+2\int_{\boldsymbol{q}}\lambda_{qu}(t_i,t_j;(\boldsymbol{p}-\boldsymbol{q})^2)\partial_{\tau}\big(G^{R}+G^{A}\big)(t_i,t_j;\boldsymbol{q})\Bigg)\,.\label{eq:qq-final}
\end{align}
Thus every two-point insertion evaluates a dressing either at vanishing channel momentum or at a single momentum $(\boldsymbol{p}-\boldsymbol{q})^{2}$, and at either one time (local tadpole) or two times (memory). This is the information that the four-point truncation must retain.

Near equilibrium, where $\lambda_{an}=0$, \Cref{eq:qc-final,eq:qq-final} reduce to the classical tadpole $\lambda_{cl}\times G^{K}$ and the quantum tadpole $\lambda_{qu}\times(G^{R}+G^{A})$, respectively.

%%%%%%%%%%%%%%%%%%%%%%%%%%%%%%%%%%%%%%%%%%%%%%%%%%%%%%%

\subsection{The flow equations for the four-point functions}
\label{sec:4pointflow}

The four-point functions include the classical, quantum, and anomaly parts. The two-point flows in \Cref{sec:2pointflow} only require the four dressings of \Cref{tab:dressings}, each a function of two times and one channel momentum. A simultaneous dependence on two channel momenta $(\boldsymbol{p}_{s},\boldsymbol{p}_{t})$ is therefore superfluous for the two-point sector. We close the four-point system by a channel-diagonal truncation: each dressing is extracted on its own momentum axis, with the crossed channel momenta set to zero.

\subsubsection{Channel-diagonal kinematics}

Instead of the two-dimensional configuration with independent $\boldsymbol{p}_{s}$ and $\boldsymbol{p}_{t}$, we use the two one-dimensional slices
\begin{align}
    \boldsymbol{p}_{i}=\boldsymbol{p}_{j}=\frac{\boldsymbol{P}}{2}\,,\qquad
    \boldsymbol{p}_{k}=\boldsymbol{p}_{l}=-\frac{\boldsymbol{P}}{2}\,,
    \label{eq:s-slice}
\end{align}
which realizes $s=\boldsymbol{P}^{2}$ and $t=u=0$, and
\begin{align}
    \boldsymbol{p}_{i}=\boldsymbol{p}_{k}=\frac{\boldsymbol{P}}{2}\,,\qquad
    \boldsymbol{p}_{j}=\boldsymbol{p}_{l}=-\frac{\boldsymbol{P}}{2}\,,
    \label{eq:t-slice}
\end{align}
which realizes $t=\boldsymbol{P}^{2}$ and $s=u=0$. The $s$-slice determines $\lambda_{cl}$, $\lambda_{qu}$, and $\lambda_{an}^{s}$; the $t$-slice determines $\lambda_{an}^{t}$ (and, as a consistency check, the $t$-channel projections of $\lambda_{cl}$ and $\lambda_{qu}$, which are identified with the $s$-channel dressings). The omitted mixed dependence of, \eg, $\lambda_{cl}^{s}$ on a nonzero $t$-channel momentum enters the vertex flow with relative weight $2/(N+4)$ and is never sampled by \Cref{eq:qc-final,eq:qq-final}.

\subsubsection{General bubble flow}

For the flow equation of the four-point function, the general form is given by
\begin{align}
    \partial_{\tau}\Gamma_{ijkl}^{(4)}(t_i,t_j,t_k,t_l;\boldsymbol{p_{i}},\boldsymbol{p_{j}},\boldsymbol{p_{k}},\boldsymbol{p_{l}})=&\nonumber\\[2ex]
    &\hspace{-5cm}\int_{\boldsymbol{q},t}\mathrm{tr}\left[\Gamma_{ii_{1}jj_{1}}^{(4)}(t_i,t_{i_1},t_j,t_{j_1};\boldsymbol{p_{i}},-\boldsymbol{q},\boldsymbol{p_{j}},\boldsymbol{q}-\boldsymbol{p_{s}})\partial_{\tau}G_{i_{1}k_{1}}(t_{i_1},t_{k_1};\boldsymbol{q})\right.\nonumber\\[2ex]&\hspace{-4cm}\times\left.\Gamma_{k_{1}kl_{1}l}^{(4)}(t_{k_1},t_k,t_{l_1},t_l;\boldsymbol{q},\boldsymbol{p_{k}},\boldsymbol{p_s}-\boldsymbol{q},\boldsymbol{p_{l}})G_{l_{1}j_{1}}(t_{l_1},t_{j_1};\boldsymbol{q}-\boldsymbol{p_s})\right]\nonumber\\[2ex]
    &\hspace{-5.3cm}+\int_{\boldsymbol{q},t}\mathrm{tr}\left[\Gamma_{ii_{1}jj_{1}}^{(4)}(t_i,t_{i_1},t_j,t_{j_1};\boldsymbol{p_{i}},-\boldsymbol{q}-\boldsymbol{p_{s}},\boldsymbol{p_{j}},\boldsymbol{q})G_{i_{1}k_{1}}(t_{i_1},t_{k_1};\boldsymbol{q}+\boldsymbol{p_s})\right.\nonumber\\[2ex]&\hspace{-4cm}\times\left.\Gamma_{k_{1}kl_{1}l}^{(4)}(t_{k_1},t_k,t_{l_1},t_l;-\boldsymbol{q},\boldsymbol{p_{k}},\boldsymbol{p_s}+\boldsymbol{q},\boldsymbol{p_{l}})\partial_{\tau}G_{l_{1}j_{1}}(t_{l_1},t_{j_1};\boldsymbol{q})\right]\nonumber\\[2ex]
    &\hspace{-5.3cm}+\int_{\boldsymbol{q},t}\mathrm{tr}\left[\Gamma_{ii_{1}kk_{1}}^{(4)}(t_i,t_{i_1},t_k,t_{k_1};\boldsymbol{p_{i}},-\boldsymbol{q},\boldsymbol{p_{k}},\boldsymbol{q}-\boldsymbol{p_{t}})\partial_{\tau}G_{i_{1}j_{1}}(t_{i_1},t_{j_1};\boldsymbol{q})\right.\nonumber\\[2ex]&\hspace{-4cm}\times\left.\Gamma_{j_{1}jl_{1}l}^{(4)}(t_{j_1},t_j,t_{l_1},t_l;\boldsymbol{q},\boldsymbol{p_{j}},\boldsymbol{p_t}-\boldsymbol{q},\boldsymbol{p_{l}})G_{l_{1}k_{1}}(t_{l_1},t_{k_1};\boldsymbol{q}-\boldsymbol{p}_{t})\right]\nonumber\\[2ex]
    &\hspace{-5.3cm}+\int_{\boldsymbol{q},t}\mathrm{tr}\left[\Gamma_{ii_{1}kk_{1}}^{(4)}(t_i,t_{i_1},t_k,t_{k_1};\boldsymbol{p_{i}},-\boldsymbol{q}-\boldsymbol{p_{t}},\boldsymbol{p_{k}},\boldsymbol{q})G_{i_{1}j_{1}}(t_{i_1},t_{j_1};\boldsymbol{q}+\boldsymbol{p}_{t})\right.\nonumber\\[2ex]&\hspace{-4cm}\times\left.\Gamma_{j_{1}jl_{1}l}^{(4)}(t_{j_1},t_j,t_{l_1},t_l;-\boldsymbol{q},\boldsymbol{p_{j}},\boldsymbol{p_t}+\boldsymbol{q},\boldsymbol{p_{l}})\partial_{\tau}G_{l_{1}k_{1}}(t_{l_1},t_{k_1};\boldsymbol{q})\right]\nonumber\\[2ex]
    &\hspace{-5.3cm}+2\int_{\boldsymbol{q},t}\mathrm{tr}\left[\Gamma_{ii_{1}ll_{1}}^{(4)}(t_i,t_{i_1},t_l,t_{l_1};\boldsymbol{p_{i}},-\boldsymbol{q},\boldsymbol{p_{l}},\boldsymbol{q})\partial_{\tau}G_{i_{1}j_{1}}(t_{i_1},t_{j_1};\boldsymbol{q})\right.\nonumber\\[2ex]&\hspace{-4cm}\times\left.\Gamma_{j_{1}jk_{1}k}^{(4)}(t_{j_1},t_j,t_{k_1},t_k;\boldsymbol{q},\boldsymbol{p_{j}},-\boldsymbol{q},\boldsymbol{p_{k}})G_{k_{1}l_{1}}(t_{k_1},t_{l_1};\boldsymbol{q})\right]\,.\label{eq:4point-flow}
\end{align}
On the $s$-slice \Cref{eq:s-slice} one has $\boldsymbol{p}_{s}=\boldsymbol{P}$ and $\boldsymbol{p}_{t}=0$; on the $t$-slice \Cref{eq:t-slice} one has $\boldsymbol{p}_{t}=\boldsymbol{P}$ and $\boldsymbol{p}_{s}=0$. The last (weight-$2$) diagram is independent of the channel momentum and is evaluated at vanishing momentum.

With the projection operators of the classical and anomaly vertices in $s$ and $t$-channels, which are given by
\begin{align}
    \left(P_{cl}^{s}\right)_{ijkl}&=\frac{3}{N(N-1)(N+2)}\left(\delta_{il}\delta_{jk}+\delta_{ik}\delta_{jl}-(1+N)\delta_{ij}\delta_{kl}\right)\nonumber\\[2ex]
    \left(P_{cl}^{t}\right)_{ijkl}&=\frac{3}{N(N-1)(N+2)}\left(\delta_{il}\delta_{jk}-(1+N)\delta_{ik}\delta_{jl}+\delta_{ij}\delta_{kl}\right)\nonumber\\[2ex]
    \left(P_{an}^{s}\right)_{ijkl}&=\frac{\mathrm{i}}{3}\left(P_{cl}^{s}\right)_{ijkl}\,,\quad\left(P_{an}^{t}\right)_{ijkl}=\frac{\mathrm{i}}{3}\left(P_{cl}^{t}\right)_{ijkl}\,,
    \label{eq:projectors}
\end{align}
we obtain the flows of the dressings. In particular, for the four-point vertices in the flow in \Cref{eq:4point-flow}, we need to exchange the internal indices together with the corresponding times and momenta, such that their definitions match those in \Cref{eq:classical-vertex} and \Cref{eq:anomaly-vertex}.

\subsubsection{Reduced dressing flows}

After projection, the self-channel bubble (flavor trace $N+4$) carries the external channel momentum $\boldsymbol{P}$, while the crossed-channel bubbles (relative weight $4$) sit at vanishing momentum. Suppressing Keldysh labels, the classical dressing on the $s$-slice obeys
\begin{align}
    \partial_{\tau}\lambda_{cl}(t_a,t_b;\boldsymbol{P}^{2})
    &=(N+4)\,\tilde{\partial}_{\tau}\!\int_{t_{1},t_{2}}\lambda_{cl}(t_a,t_{1};\boldsymbol{P}^{2})\,\Pi(t_{1},t_{2};\boldsymbol{P})\,\lambda_{cl}(t_{2},t_b;\boldsymbol{P}^{2})\nonumber\\[2ex]
    &\quad +4\,\tilde{\partial}_{\tau}\!\int_{t_{1},t_{2}}\lambda_{cl}(t_a,t_{1};0)\,\Pi(t_{1},t_{2};0)\,\lambda_{cl}(t_{2},t_b;0)\,,
    \label{eq:lambda-cl-flow}
\end{align}
and $\lambda_{qu}$ satisfies the same structure with the propagator combination appropriate to $qqqc$. Here $\Pi$ denotes the one-loop bubble built from the two-point propagators in \Cref{eq:4point-flow}. The first line is the memory-carrying self-channel resummation; the second line is a local renormalization of the coupling and may be evaluated at equal time, $t_{1}=t_{2}$.

The anomaly dressings are projected separately. The $s$-channel dressing $\lambda_{an}^{s}$ is extracted with $P_{an}^{s}$ on \Cref{eq:s-slice}; the $t$-channel dressing $\lambda_{an}^{t}$ is extracted with $P_{an}^{t}$ on \Cref{eq:t-slice}. Their flows have the same split into a self-channel two-time convolution at momentum $\boldsymbol{P}$ and a weight-$4$ equal-time piece at zero momentum, but they are not identified with each other.

\subsubsection{Time truncation}

The channel ansatz already collapses four external times onto the two endpoint times of each channel. The remaining cost is the two-time convolution in \Cref{eq:lambda-cl-flow}. We keep this two-time dependence in the self-channel bubble, since it is exactly the memory $t_i-t_j\neq 0$ required by \Cref{eq:qc-final,eq:qq-final}. Two optional further reductions, not used in the baseline truncation, are:
\begin{itemize}
\item a derivative expansion in the center time $T=(t_a+t_b)/2$, retaining the full relative-time dependence $\Delta t=t_a-t_b$ at zeroth order in $\partial_{T}$;
\item replacing the weight-$4$ crossed bubble by its equal-time, zero-momentum value, as already indicated in \Cref{eq:lambda-cl-flow}.
\end{itemize}

\subsubsection{Anomaly sector}

Since $\lambda_{an}$ vanishes in equilibrium, a practical hierarchy is to first solve the coupled system of $\lambda_{cl}$, $\lambda_{qu}$ and the two-point functions with $\lambda_{an}=0$, and only afterwards generate $\lambda_{an}^{s}$ and $\lambda_{an}^{t}$ from the right-hand side of \Cref{eq:4point-flow}. Full backcoupling of the anomaly dressings is required only if they grow to a size comparable with $\lambda_{cl}$.


\hfill
\bibliography{ref-lib}% Produces the bibliography via BibTeX.
%%%%%%%%%%%%%%%%%%%%%%%%%%%%%%%%%%%%%%%%%%%%%%%%%%%%%%%%%%

\end{document}
